\documentclass{article}

\usepackage{amsmath,amssymb}
\usepackage{xurl}
\usepackage[hidelinks]{hyperref}
\setlength{\emergencystretch}{2em}

\input{command}

\title{Wolf's PPT--Indecomposable Existence Equivalence:\\
Resolution of the Reverse Implication}
\author{}
\date{2026-08-24}

\begin{document}
\maketitle
\gapnote{scope-restriction}{resolved}

\section{The source equivalence}

For fixed dimensions $d,d'$, Wolf's Proposition~3.5 states
\begin{align}
  &\text{there is an indecomposable positive map }
      T:\MN{d}\longrightarrow\MN{d'}
  \notag\\
  &\qquad\Longleftrightarrow\qquad
  \begin{gathered}
    \text{there is an entangled density operator }\rho
      \in\MN{d}\otimes\MN{d'}\\
    \text{with }\rho^{T_1}\geq 0.
  \end{gathered}
  \tag{Wolf Proposition 3.5}
\end{align}
Immediately before the proposition, Wolf writes the decomposable witnesses as
\begin{align}
  W&=P_1+P_2^{T_1},
  &P_1,P_2&\geq0.
  \tag{Wolf 3.15}
\end{align}
Thus the transpose in both the witness decomposition and the PPT condition is
on the first factor, of dimension $d$. The local source is
\path{Notes/WolfNoteTexSource/ch03_positive_not_completely.tex}, lines
276--303.\footnote{Tracked by \ghissue{22}.}

The witness argument used below is the normalized form of Theorem~3 of
Lewenstein--Kraus--Cirac--Horodecki
\cite{LewensteinKrausCiracHorodecki2000}. With trace-one positive operators
$P_1,P_2$, its compact convex set, written in Wolf's first-factor convention,
is
\begin{align}
  \mathcal D_1
  &=\left\{aP_1+(1-a)P_2^{T_1}:
      0\leq a\leq1,\ P_1,P_2\geq0,\
      \tr(P_1)=\tr(P_2)=1\right\}.
  \tag{normalized decomposable set}
\end{align}

\section{The rectangular formal statement}

For $d,d'\geq1$, the formalized equivalence has the same input-to-output
orientation as Wolf:
\begin{align}
  &\exists T:\MN{d}\longrightarrow\MN{d'}
      \text{ positive and indecomposable}
  \notag\\
  &\qquad\Longleftrightarrow\qquad
  \begin{gathered}
    \exists\rho\in\MN{d}\otimes\MN{d'}:\
      \rho\geq0,\quad\tr(\rho)=1,\\
      \rho^{T_1}\geq0,\quad\rho\text{ not separable}.
  \end{gathered}
  \tag{formalized Wolf Proposition 3.5}
\end{align}
The fixed nonzero-dimensional theorem is
\leanid{Matrix.exists_isIndecomposablePositiveMap_iff_exists_isPPT_not_isSeparable}.
The theorem
\leanid{Matrix.exists_isIndecomposablePositiveMap_iff_exists_isPPT_not_isSeparable_all_dimensions}
states the same equivalence for arbitrary natural-number dimensions and
settles the zero-dimensional boundary explicitly.\footnote{Formal module:
\doclink{QICLean/Channel/PPTIndecomposable}.}

The implication from a PPT entangled state to an indecomposable map is
\leanid{Matrix.exists_isIndecomposablePositiveMap_of_isPPT_not_isSeparable}.
It applies Wolf's rectangular positive-map criterion at level one. If the
detecting map were decomposable, writing $T=T_1+T_2\theta$ would give
\begin{align}
  (T\otimes\operatorname{id})(\rho)
    &=(T_1\otimes\operatorname{id})(\rho)
      +(T_2\otimes\operatorname{id})(\rho^{T_1})\geq0,
  \notag
\end{align}
contrary to detection.

\section{The reverse implication}

Let $T:\MN{d}\to\MN{d'}$ be indecomposable and positive, with
$d,d'\geq1$. Let $T^*:\MN{d'}\to\MN{d}$ denote its trace-pairing adjoint,
and define
\begin{align}
  W_0&=\tau(T^*)\in\MN{d}\otimes\MN{d'}.
  \tag{trace-adjoint Choi witness}
\end{align}
This is the output-first rectangular Choi matrix of $T^*$. Its rows and
columns are indexed by
$\operatorname{Fin}(d)\times\operatorname{Fin}(d')$, in that order. Hence
the first factor of $W_0$ has dimension $d$, exactly as in Wolf's
Equation~(3.15); no exchange of the tensor factors occurs.

The trace adjoint of a positive map is positive. The rectangular Choi
criterion at level one therefore gives
\begin{align}
  \operatorname{Re}\tr\!\left(
    W_0\ketbra{\psi}{\psi}\right)&\geq0
  \qquad
  \text{if $\psi$ has Schmidt rank at most one}.
  \tag{block positivity}
\end{align}
Moreover, indecomposability implies $T\neq0$. Positivity shows that
$T(\Id_d)=0$ would force $T=0$, so $T(\Id_d)$ is a nonzero positive
operator. The rectangular Choi trace identity gives
\begin{align}
  \tr(W_0)&=\frac{1}{d'}\tr\!\left(T(\Id_d)\right)>0.
  \tag{positive witness trace}
\end{align}
Consequently
\begin{align}
  W&=\frac{W_0}{\tr(W_0)}
  \tag{normalized witness}
\end{align}
is Hermitian, block-positive, and trace one. The trace-adjoint rectangular
Choi correspondence identifies decomposability of $T$ with membership of
$W_0$ in Wolf's cone $P_1+P_2^{T_1}$. Positive rescaling does not change
that membership, so $W\notin\mathcal D_1$.

Strictly separate $W$ from the compact convex set $\mathcal D_1$ in the
real space of complex matrices. Represent the separating real-linear
functional by a Hermitian matrix and subtract a multiple of the trace
functional. Since $W$ and every member of $\mathcal D_1$ have trace one,
this gives a Hermitian operator $R$ such that
\begin{align}
  \operatorname{Re}\tr(WR)&<0,
  &\operatorname{Re}\tr(DR)&\geq0
    \qquad(D\in\mathcal D_1).
  \tag{normalized separation}
\end{align}
At the endpoint $a=1$, the second inequality holds for every positive
trace-one $P$. Scaling then extends it to every $P\geq0$, and self-duality
of the positive cone gives $R\geq0$. At the endpoint $a=0$, it holds for
$Q^{T_1}$, where $Q$ is positive and has trace one. The identity
\begin{align}
  \operatorname{Re}\tr\!\left(Q^{T_1}R\right)
    &=\operatorname{Re}\tr\!\left(QR^{T_1}\right)
  \tag{first-factor transpose adjointness}
\end{align}
gives $R^{T_1}\geq0$. This is the first-factor transpose throughout.

The strict negative pairing implies $R\neq0$. Hence $\tr(R)>0$, and
\begin{align}
  \rho&=\frac{R}{\tr(R)}
  \tag{PPT density operator}
\end{align}
is positive, has trace one, and satisfies $\rho^{T_1}\geq0$. Finally,
$\operatorname{Re}\tr(W\rho)<0$, whereas block positivity of $W$ makes its
pairing nonnegative on every separable operator. Thus $\rho$ is entangled.

The separation and normalization steps are formalized by
\leanid{Matrix.exists_posSemidef_isPPT_negative_separator_of_not_isDecomposableWitness}
and
\leanid{Matrix.exists_isPPT_not_isSeparable_of_not_isDecomposableWitness}.
The complete map-to-state implication is
\leanid{Matrix.IsIndecomposablePositiveMap.exists_isPPT_not_isSeparable}.
The compactness and convexity of $\mathcal D_1$, the first-factor trace
identity, and the rectangular Choi correspondence are respectively
\leanid{Matrix.isCompact_setOf_isNormalizedDecomposableWitness},
\leanid{Matrix.convex_setOf_isNormalizedDecomposableWitness},
\leanid{Matrix.trace_partialTransposeLeft_mul_re}, and
\leanid{ChoiRectangular.isDecomposablePositiveMap_iff_choiMatrix_traceAdjointMap_isDecomposableWitness}.

This proof uses neither closedness of the full unbounded cone nor a
definition of decomposability by duality. It separates from the compact
trace-one set and proves the two positivity conclusions from its endpoints,
as in the normalized witness argument of
Lewenstein--Kraus--Cirac--Horodecki.

\section{The zero-dimensional boundary}

If $d=0$, every map with domain $\MN{0}$ is the zero map. If $d'=0$, every
map with codomain $\MN{0}$ is again the zero map. In either case the map is
decomposable, so no indecomposable positive map exists. At the same time,
the tensor product is empty and every matrix on it has trace zero, so no
density operator exists. Both sides of the equivalence are therefore false.
This proves the arbitrary-dimensional statement without imposing a hidden
nonzero-dimensional hypothesis.

\section{Verdict}

The scope restriction recorded here is resolved. Both directions of Wolf
Proposition~3.5 are formalized with independent dimensions $d,d'$, the
printed map orientation $\MN{d}\to\MN{d'}$, the trace-adjoint rectangular
Choi witness on $\C^d\otimes\C^{d'}$, and partial transpose on the first
factor. The fixed nonzero-dimensional equivalence and the explicit
zero-dimensional boundary are both proved. No source-level gap remains in
this proposition.

\bibliographystyle{alpha}
\bibliography{references}

\end{document}
